\section*{The Propers: Detailed Commentary}

\subsection*{Solomon asks to hear before he judges\enspace\elementref{First Reading; Collect}}

The lectionary preserves Solomon's confession of youth and inexperience, the people entrusted to him, the request for a discerning heart, and God's approval. Verse 6, which the cut omits, is not part of the proclaimed text. The request is public in consequence even though received in a dream: wisdom is sought for just governance, not private distinction.

Ambrose cites the requested prudent heart in \emph{On the Duties of the Clergy} II.8, 44--49 as a heart able to hear and judge justly, then develops the later judgment between two mothers. He illuminates the union of prudence and justice; the appointed pericope itself, not Ambrose, supplies the contrast with riches and longevity.

Benedict XVI's Angelus on this exact Sunday in 2011 moves from Solomon's hearing heart to conscience: listening to truth enables the distinction of good and evil and places governing responsibility under justice and peace. Conscience is Benedict's word for that heart, not an alternate translation of 1~Kings; his address sharpens the reading's public consequence without reducing biblical heart-language to one later technical term.

The broader canonical story tempers the scene. The king who asks well can later accumulate wealth, wives, and divided allegiance. Wisdom received is not fidelity guaranteed. The Collect's repeated dependence---God as protector, ruler, and guide---keeps discernment from becoming self-certifying expertise.

Giuseppe Ferraro's biblical commentary on this Collect in \emph{Notitiae} traces protector to Vulgate Psalm 17, dependence to the vine discourse of John 15, ruler and guide to biblical governance and Exodus, and the passing/enduring contrast to 1~Corinthians 7. Ferraro explicitly refuses the inference that detachment cancels work within the world: created undertakings receive durable value in communion with Christ. The article is a signed study in the Dicastery's journal, not a magisterial ruling.

Peter Comestor's medieval \emph{Historia scholastica} retells the dream while preserving the distinction between Solomon's request and God's additional gifts. It witnesses to the episode's long pedagogical life but is not an independent source for the ancient event. Later historical and devotional commentators likewise admire the request while the canonical sequel controls idealization.

\begin{threestable}{Movement}{Meaning}{Limit}
\textbf{Confess inadequacy} & Office does not create the wisdom it requires. & Humility does not excuse refusal to learn or decide.\\
\textbf{Ask for a hearing heart} & Discernment serves justice for a concrete people. & It is not intuition exempt from evidence, counsel, or law.\\
\textbf{Receive gift} & God approves a request ordered beyond advantage. & Gift does not make later choices infallible.\\
\end{threestable}

\sectionguard
\subsection*{Gathered strength and gifts returned\enspace\elementref{Entrance; Prayer over Offerings}}

The Entrance adapts clauses from Psalm 67 rather than proclaiming a continuous span. The whole processional hymn praises God's coming, defense of the vulnerable, gift of a household, and strength for the people. Augustine treats divine presence, gathering, and strength as gifts rather than human possession (\emph{Exposition on Psalm 68}, 6, 8, 40); Bellarmine later receives the same movement within ecclesial praise.

The Prayer over the Offerings asks God to receive gifts that first come from divine bounty and to make the mysteries sanctify present conduct. The ritual movement is reception, return, and transformation: the assembly does not purchase grace with its own property. This shared Week XVII prayer may illuminate the Gospel's relinquishment. There is no Missal Offertory antiphon, and the actual music used remains unknown.

\sectionguard
\subsection*{Instruction, treasure, and pearl\enspace\elementref{Responsorial Psalm; Gospel branches}}

Psalm 119 repeatedly turns value into prayer: the Lord is the singer's portion; instruction exceeds gold and silver; commands are loved above fine gold; their unfolding gives light. The refrain's love of the law is not abstract admiration but a sustained appetite formed by meditation, mercy, and obedience.

Ambrose's \emph{Exposition of Psalm 118} VIII, 2--4 asks how rare it is to say that the Lord is one's portion without reserving mastery to wealth or appetite. In IX, 21--22 he treats the law above gold and silver as preference proved in action and relates the words of God's mouth to words of eternal life. Augustine keeps the appointed loci's portion, love of law, covenant mercy, wonder, and understanding together; Bellarmine likewise connects valuation and the request for light. None of them expounds the later selection as a unit; its separated stanzas answer to verse 97a, and the Psalm's Israelite Torah meditation remains the literal context for their Christian reading.

The two parables intensify valuation. One person encounters treasure and acts in joy; the merchant searches and recognizes one pearl of great price. Chrysostom's \emph{Homily 47 on Matthew} emphasizes the Kingdom's hidden worth and the single-minded response it evokes. The differences remain: discovery and search, field and market, treasure and pearl are not one psychology.

Francis's Angelus for this exact Sunday in 2014 also preserves the difference between unexpected discovery and sustained search, then reads encounter with Christ as reorienting an entire life and giving meaning even to toil, suffering, and death. His 2020 Angelus holds grace and human willingness together: the Kingdom is not achieved by price, yet grace calls for an active welcome. These are contemporary homilies on this Sunday's readings.

Joy does not purchase the Kingdom. In both images the incomparable discovery reorganizes possession; the price imagery cannot make salvation a commodity. Nor can ``selling all'' be abstracted into irresponsible disposal of dependents' goods or contempt for material creation. The Psalm itself values instruction through a material comparison without calling gold evil.

\sectionguard
\subsection*{The long branch: net, judgment, and a scribe with both treasures\enspace\elementref{long Gospel}}

Only the long form proceeds to the net gathering every kind, the shore-side sorting, the angels' final separation, and the householder's new and old treasure. The net gives eschatological seriousness to value: present gathering does not erase final judgment. It is not a mandate for disciples to conduct the angelic sorting now.

Origen's \emph{Commentary on Matthew} X treats the trained scribe and the relation of new and old learning. Augustine's \emph{Sermon 24} likewise insists that the instructed teacher brings both forth rather than discarding one.

Aquinas's own \emph{Super Matthaeum} XIII, lect.~4 arranges the closing parables as dimensions of Gospel teaching: abundance in the treasure, beauty or charity in the pearl, and universal reach in the net, whose shore-side sorting qualifies the present gathering. At verse 52 he turns Jesus' question about understanding toward teaching office: the instructed scribe draws new and old from sacred doctrine's treasury and, as Christ's notary, writes the commands first in his own heart and then in others. Aquinas's \emph{Catena aurea} places Chrysostom's fearful net after treasure and pearl so that hearers do not rely on preaching or bare faith without a corresponding life. These are medieval syntheses; the Catena compiles patristic extracts, and Aquinas is their compiler, not their author. Matthew's own image calls both stores treasure.

The short form ends before all of this. A celebration choosing Matthew 13:44--46 does not proclaim net, angels, furnace, understanding, or the scribe. These belong to the long form, an authorized alternative, and preaching claims must follow the branch actually proclaimed.

\sectionguard
\subsection*{Called toward the Son amid groaning\enspace\elementref{Second Reading}}

Romans 8:28 is often detached from its argumentative world. Immediately before, creation and believers groan and the Spirit aids weak prayer; immediately after, Paul names affliction, persecution, hunger, peril, and the love from which none can separate the baptized. The ``all things'' claim is therefore consolation within suffering, not denial of it.

Chrysostom's \emph{Homily 15 on Romans} interprets God's action as turning even adversity toward the good of those who love him. He does not call the adversity good in itself. The chain of foreknowledge, predestination, call, justification, and glorification is centered on conformity to the Son, who is firstborn among many brothers and sisters. It supports confidence in divine initiative, not fatalism about individual events or a prosperity guarantee.

Aquinas first distinguishes the good of the whole from the good of the person who suffers, then reads Romans 8:28 as the Spirit's particular ordering for those who love God: affliction can exercise humility, patience, wisdom, and benevolence without becoming good in itself. At verses 29--30 he treats conformity to the Son as the end or effect of predestination, not a foreseen merit that causes it, and unfolds the chain from eternal preparation into calling and justification (\emph{Super Romanos} VIII, lect.~6).

The Second Vatican Council places this conformity in a paschal horizon: Christ makes the human vocation clear, and the person patterned on the dying Christ advances toward resurrection in hope (\emph{Gaudium et spes}, 22). The \emph{Catechism} supplies an equally important boundary at nn.~311--312: creaturely freedom can produce moral evil, God is not its direct or indirect cause, and the good providence draws from evil never makes evil itself good. Together these sources define both the confidence and the refusal required by Romans 8.

Haydock's later compilation makes complementary moves. At verse 28, ``purpose'' is God's will rather than a human plan. At verse 29, conformity includes suffering with Christ, following his teaching, and imitating his life; the notes join special grace to free cooperation and warn against prying into secret election. At verse 30 they name faith, hope, love, and penance in justification and do not infer final perseverance from every justified state. Divine initiative therefore does not erase agency or authorize assigning a living person a certain final status.

The reading is semi-continuous. Its ``purpose'' belongs to Paul's argument, not to Solomon's administrative vocation or the merchant's search. Read side by side, the readings can ask what sort of good governs discernment, but Matthew's treasure does not decode Paul's predestination.

\begin{studybox}{Trust that does not bless evil}
The reading does not tell a harmed person that abuse, catastrophe, or injustice was secretly good, nor that protest or protection resists Providence. Christian hope says God can work within what creatures do wickedly and conform sufferers to Christ without blessing evil. Immediate safety, lament, justice, medicine, and communal aid remain compatible with trust.
\end{studybox}

\sectionguard
\subsection*{Revelation received by the humble\enspace\elementref{Acclamation}}

The acclamation adapts rather than reproduces Jesus' thanksgiving in Matthew 11:25. Chrysostom's \emph{Homily 38} says the contrast does not condemn genuine learning; it exposes self-sufficient resistance and praises those who receive revelation without pretension. Aquinas likewise distinguishes the proud posture from teachable littleness. That humility fits Solomon's request but is not identical with royal discernment, and Matthew 11 belongs to a different scene from Matthew 13.

\sectionguard
\subsection*{Two Communion paths and one final petition\enspace\elementref{Communion A; Communion B; Prayer after Communion}}

Communion A takes one verse from Psalm 102 under the Latin numbering: the soul commands itself to bless and not forget divine benefits. Augustine's exposition of Psalm 103 (English page numbering) makes the command self-address and unfolds memory through mercy, healing, rescue, compassion, and creaturely frailty. Bellarmine likewise treats the opening summons as deliberate gratitude. At Communion that canonical memory resists entitlement and amnesia, but neither commentary proves an original Eucharistic referent.

Communion B joins two adjacent beatitudes without making them one merit formula. Augustine distinguishes merciful action and its promised mercy from the purified heart and its promised vision; Chrysostom emphasizes generosity and an interior freedom from corrupt desire, while Aquinas later keeps the promised ends distinct. None comments on this later antiphon. At the Eucharist the pair places reception under a demanding horizon: the communicant asks to become merciful and undivided, not merely to possess a sacred object.

The alternatives are exclusive, and either can illuminate the final prayer's confession of the perpetual memorial of the Son's Passion and its petition that divine love advance salvation. Its fruit is gift calling for lived response, neither mechanical nor self-achieved.
